\begin{exercise}[The Pontryagin density as an exact form]{I-CH10-013}
In $D=4$ Euclidean dimensions, with orientation $\epsilon^{0123}=+1$, define
$\widetilde F^{\mu\nu}=\tfrac12\epsilon^{\mu\nu\rho\sigma}F_{\rho\sigma}$.
First take an Abelian potential $A$.  Use $F=\dd A$ and $\dd F=0$ to show
\[
  F\wedge F=\dd(A\wedge F),
\]
and give the corresponding component total derivative for
$F_{\mu\nu}\widetilde F^{\mu\nu}$.

Next use Banks's anti-Hermitian matrix-valued conventions
\[
  A=\ii A_\mu\dd x^\mu,\qquad F=\dd A-A^2,
  \qquad C=\operatorname{tr}\left(A\,\dd A-\frac23A^3\right).
\]
Keeping the graded signs visible, expand $\dd C$ and
$\operatorname{tr}F^2$ and prove
$\dd C=\operatorname{tr}F^2$.  State the relation between
$\operatorname{tr}F\wedge F$ and
% CONVENTION-SCOPE: fixed D=4 reason=Pontryagin density in D=4 Euclidean space
$\operatorname{tr}(F_{\mu\nu}\widetilde F^{\mu\nu})\dd^4x$ under the
convention $F=\tfrac12F_{\mu\nu}\dd x^\mu\wedge\dd x^\nu$.
\end{exercise}
